%! Author = Saeid Yazdani
%! Date = 7/18/2024


\section{Time-Variant Power Supply Control Loop}
\label{subsec:time-variant-comp-network}

A pure analog control loop is classified as a \emph{Time-Invariant} system
since the output response of the system remains consistent regardless
of the point of time an input is applied.
The system's behavior and characteristics are fixed, determined by its physical
components and their configurations.
In such systems, for any given input, the output will be the same, except due to
ageing-related drift of the components or influence of noise in the environment.

On the other hand, a hybrid digital-analog control loop represents
a \emph{Time-Variant} system.
The digital section introduces adaptability and real-time response.
The digital components can modify their behavior based on pre-programmed
algorithms or sets of conditions, allowing the system to adjust its output
dynamically in response to changes in the input.
This ability to alter the output on the fly provides flexibility and precision
in applications where the operating conditions can vary widely or require rapid
adjustments.

The primary challenge in implementing adaptive compensation in analog circuitry
lies in the switch controller, which is digital rather than analog as
shown in \autoref{fig:traditional-psu-loop}.
This inherent digital nature necessitates the integration of a switching
processor, enabling signal processing in the digital domain through a
series of \ac{ADC} and \ac{DAC} conversions.

\begin{figure}
    \centering
    \printfitgraphics[0.58\textheight]{figures/dps_algorithms_psu/traditional-analog-power-supply-circuit}
    \caption[Comparison of Time-Invariant and Time-Variant ]{
        Traditional analog power supply control loop.
        \emph{Power Stage} Converts the input power to the desired output power
        and typically includes main power switch (\acs{MOSFET} or \acs{BJT})
        controlled by the duty cycle of a PWM generator.
        The \emph{Feedback Network} senses the output voltage or current and
        feeds back to the error amplifier for comparison with a reference
        voltage.
        Stabilization of the control loop is done by \acs{RLC}-based filters in
        the \emph{Compensation Network}.
    }
    \label{fig:traditional-psu-loop}
\end{figure}

Adaptive compensation can be more effectively achieved by utilizing \ac{DSP}.
DSPs allow for precise and flexible adjustments, accommodating a wide range of
operating conditions and enhancing overall system performance.
Furthermore, DSP-based systems facilitate real-time monitoring and adjustment
of compensation parameters, providing an optimal balance between stability and
response time.
This dynamic capability ensures that the system can efficiently respond to
changes in load and operating conditions, maintaining performance and preventing
issues such as overshoot or instability.

In the proposed power supply we use \ac{DSP} for adaptive control algorithms,
on-the-fly compensation, and signal monitoring for safety reasons
(e.g. clamping functions).

\autoref{fig:invariant-vs-variant-control} demonstrates the dynamic adjustment
of compensation networks to slow down the amplifiers just before overshoot
occurs.
This technique requires parameter sets that are dependent on both the range
and the load conditions.
The necessary parameters are stored in a memory organized by the processor,
with all loop calculations executed by a high-speed \ac{DSP} core within
a \ac{FPGA}.
This method ensures precise control and optimal performance by dynamically
adapting to varying operational conditions.

\begin{figure}[htbp]
    \centering
    \printfitgraphics[0.62\textheight]{figures/dps_algorithms_psu/time-invariant_vs_time-variant_control}
    \caption[Comparison of Time-Invariant and Time-Variant ]{
        Predictive and adaptive compensation settings in a digital control loop.
        Pure analog compensation (top) and adaptive compensation (bottom) where
        dotted red and blue waveforms are analog and solid lines are digitally
        compensated waveform.
        When used in an \ac{ATE} environment, the \ac{DUT} characteristics
        and overall electrical characteristics of the whole setup
        (including load board data) can be fed to the control algorithm
        and be used for optimized reactions to changes in the load conditions.

    }
    \label{fig:invariant-vs-variant-control}
\end{figure}

The advantages of \ac{DSP}-based digital control are based on programmability
with \emph{Time-Variant} parameters – coming from memory
\autocite{yazdani-digital-control}.

Whereas we need much instrumentation in the analog case to modify
loop parameters on the fly (multiplexers to select one out of many correction
values) we have an infinite space of parameters in the digital domain,
assuming that there are no memory constraints.

In the era of \ac{FPGA} chips with \SI{}{\giga\hertz} clock rates and
a-to-d or d-to-a converters with bandwidths of hundreds of \SI{}{\mega\hertz},
we do not suffer from speed limitations.

\autoref{fig:system-block-diagram-hvs} shows the block diagram of the new power
source for \ac{ATE} \ac{DUT} applications:

\begin{itemize}[topsep=8pt,itemsep=4pt,partopsep=4pt, parsep=4pt]
    \item
    All measurement signals are transferred to the digital domain.
    We can transfer data from high potentials to ground level with perfect
    data couplers.
    Isolation barriers of many \SI{}{\kilo\volt} are no problem, and there is no
    degradation of signal quality as we transfer digital signals only.
    \begin{itemize}
        \item
        A pure analog solution has drawbacks: Isolation amplifiers have a
        low bandwidth that lowers overall speed, and resistive dividers suffer
        significantly from accuracy problems.
    \end{itemize}

    \item
    All amplifier or switching stages in the signal-processing chain can be
    directly handled by the loop controller.
    Therefore, it is possible to
    inject compensation signals into every stage to optimize the
    output waveform shape.

\end{itemize}

\begin{figure}[p]
    \centering
    \printfitgraphics[0.80\textheight]{figures/dps_algorithms_psu/digital-feedback-structure}
    \caption[System Overview of Time-Variant Control Loop]{
        System overview of the proposed Time-Variant Control loop for the
        improved \ac{DPS}, integrating both analog and digital components.
        The analog domain includes the main, voltage, and current amplifiers
        with local and global compensation to maintain stable and accurate
        output.
        The digital domain utilizes \ac{ADC} and \ac{DAC} for information
        exchange between the two domains, an \ac{MCU} for
        mode control and safety, and an FPGA for executing the control
        algorithms.
        The system is designed to manage set points, compensate for delays
        and phase shifts to provide precise output for various types
        of \acp{DUT}.
    }
    \label{fig:system-block-diagram-hvs}
\end{figure}

%\newpage
The benefits of such an arrangement are:
\begin{itemize}[topsep=8pt,itemsep=4pt,partopsep=4pt, parsep=4pt]
    \item
    The feedback controller can apply predictive algorithms
    \item
    Range and load-dependent settings are part of the control logic (\ac{FPGA})
    – we can store many parameter sets as required.
    \item
    Feedback parameters can be modified on the fly even during the settling
    of the amplifiers.
\end{itemize}

